\scanpage{194}
% Large diagonal strokes cancel both paragraphs. The last line is struck separately.
\deleted{\uline{Proposition} Supposons R1 R3 R4. Alors R2 équivaut à la condition que $r$ est \uline{fidèle}.}

\deleted{Supposons $r$ fidèle, prouvons R2. Soient $X'$, $X''$ deux ss-objets de $X\in\mathrm{Ob}\,\mathscr R_{S'}$ tels que $r(X')=r(X'')$ comme sous-objet de $r(X)$ ; \deleted{donc $r(X')\unsure{\cap}$} or $r(X'\cap X'')=r(X')\cap r(X'')$, (cf lemme) donc l'inclusion $\deleted{\unsure{r(|\text{et}}}X'\cap X''\subset X'$ et $X'\cap X''\subset X''$ \unsure{induit|induit un|deviennent des} isom sur $r$. On est ramené au cas où on a \inserted{déjà} une inclusion $X'\subset X''$, i.e. $X''=X'\amalg X_1$. Alors $r(X'')=r(X')\amalg r(X_1)$, \deleted{et l'hyp: $r(X')=r(X'')$ équivaut à $r(X_1)=\varnothing$.}}
% archive sheet number: 193
